% Методика виміряного пошуку; невиконані вимірювання ресурсів позначено явно.
\section{Методика оцінювання}\label{sec:method}

\subsection{Системи й одиниця ранжування}
Усі системи використовують одну реалізацію пошуку archive-index з явними перемикачами конфігурації. Порівнюються BM25 без лем і з лемами, повнорозмірна EmbeddingGemma~2, варіант EG2 із 256 вимірами, multilingual-E5, злиття взаємних рангів, адаптивна система та три абляції: без переранжування фраз, обробки ідентифікаторів або фільтра коротких фрагментів. Прогони мають рівень документів: ранг визначає перший знайдений фрагмент документа, наступні його фрагменти згортаються. Для запиту зберігається до 100 документів. Індекс E5 лише текстовий, на відміну від повних індексів EG2; відмінність наборів кандидатів описано як обмеження.

\subsection{Набір запитів і метрики}
Для кожної системи й варіанта суддів середні метрики обчислюються на тих самих 170 запитах із відповіддю з однаковою вагою кожного запиту. Тридцять запитів, позначених як такі, що не мають відповіді, залишаються у файлах релевантності, але не входять до цих середніх; для них окремо наводяться порожня видача й наявність результатів з позитивною оцінкою. Запит із відповіддю, для якого в пулі немає позитивних оцінок, залишається в середньому з нульовою метрикою.

Основна метрика — nDCG@10 із лінійними вагами релевантності 0, 1 і 2, дисконтуванням на $\log_2(i+1)$ для рангу $i$ й нормалізацією за ідеальним ранжуванням оцінених документів. Додаткові метрики — обернений ранг на глибині 10, успіх на глибині 1 і 5 та повнота на глибині 100. У бінарних метриках релевантними вважаються оцінки 1 і 2. Перші десять знайдених документів кожної системи оцінено повністю, проте поза об'єднаним пулом глибини 10 оцінки неповні. Отже, повнота на глибині 100 визначається відносно відомого релевантного пулу, а не всього корпусу. Результати також групуються за категорією й мовою запиту.

\subsection{Невизначеність і порівняння}
Довірчі інтервали середніх метрик і парних середніх різниць отримано процентильним бутстрепом за запитами: 10\,000 перевибірок, початкове зерно 2026. Двобічний парний рандомізаційний тест зі зміною знаків використовує 10\,000 перевибірок і поправку $(b+1)/(10\,000+1)$. Для кожної метрики й варіанта суддів дев'ять порівнянь з повною адаптивною системою утворюють одну сім'ю; p-значення коригуються методом Holm усередині сім'ї з порогом 0,05. Початкове дослідницьке питання про морфологію додатково оцінюється окремо позначеним парним порівнянням BM25 із лемами та без них за nDCG@10; його нескориговане p-значення не входить до сім'ї порівнянь з адаптивною системою. Перетин окремих довірчих інтервалів або відсутність значущої різниці не тлумачаться як еквівалентність.

Для початкової узгодженості суддів наведено лінійно зважену каппу Коена з 95\% інтервалом за 2\,000 бутстреп-перевибірками запитів. Перевибірка запитів, а не окремих пар оцінок зберігає групування всередині запиту. Ці інтервали узгодженості відокремлені від інтервалів метрик пошуку з 10\,000 перевибірками. Для кожного варіанта релевантності наведено ранги систем і зміни попарного порядку; середні, рівні з точністю $10^{-12}$, мають спільний ранг.

\subsection{Відтворюваність і рівні свідчень}
Початкові аркуші оцінок, визначення запитів, маніфест корпусу, пошукові прогони й доступні журнали викликів зафіксовано до обчислення поданого сукупного порівняння ефективності. Маніфест SHA-256 перевіряє ці входи; окремий запис підготовки пов'язує перегляд із початковими розбіжностями та порядком подання. Структуровані відповіді записуються разом з ідентифікаторами моделей, розмов і використанням токенів, коли вони доступні. Після повного перегляду підсумкові мітки замінюють лише оцінки розбіжностей, зберігаючи всі узгоджені оцінки. Таблиці й рисунки генеруються з отриманих JSON-артефактів, а не підтримуються як незалежні числові копії.

Модульні й інтеграційні перевірки встановлюють поведінку реалізації; публічний бенчмарк дає свідчення пошуку за оцінками моделей; окремі мультимодальні експерименти використовують мітки джерельних наборів даних. Ці типи свідчень не замінюють перевірку релевантності людиною чи вимірювання під час експлуатації.

\subsection{Вимірювання ресурсів}
Затримка запиту --- це час кодування запиту та пошуку фрагментів в оцінювальній обв'язці, включно з лексичним і семантичним пошуком та злиттям, із межею 400 фрагментів-кандидатів. Таймер зупиняється до згортання до документів і визначення шляхів; це не наскрізна затримка інтерфейсу. Кожна система відповідає на всі 200 запитів один раз для розігріву і ще раз під вимірюванням, в одному процесі, спершу з видимим GPU, а потім з порожньою змінною \code{CUDA\_VISIBLE\_DEVICES} (моделі вкладень тоді працюють на CPU у float32). Під час кожного виміряного проходу обв'язка записує використання CPU по ядрах із \code{/proc/stat}, секунди CPU кожного потоку процесу з \code{/proc/self/task}, зміну лічильників переривань по ядрах із \code{/proc/interrupts}, піковий резидентний розмір і пам'ять GPU, яку \code{nvidia-smi} опитує кожні 0,2~с. Машина працювала у звичайному робочому сеансі; середнє навантаження за хвилину перед кожним кроком записано в журнал вимірювань.

Вартість індексації виміряно на фіксованій випадковій вибірці корпусу (зерно 2026: 300 нотаток, 40 документів, 15 таблиць, 60 зображень і 60 записів), яку команда \code{archive} індексує в порожню базу з GPU. Для CPU взято п’яту частину за кількістю: 60 нотаток, 8 документів, 3 таблиці, 12 зображень і 12 записів. Вибірки відрізняються розміром та складом, тому співвідношення їхнього повного часу не є парною оцінкою прискорення CPU/GPU. Розміри чотирьох повних індексів бенчмарку подано такими, як їх побудовано.

Для квантування текстові фрагменти повного індексу (84\,788 векторизованих рядків текстових колекцій) кодуються заново GGUF-збірками EmbeddingGemma~2 (f16, Q8\_0, Q4\_K\_M), які обслуговує llama.cpp із середнім пулінгом і тими самими префіксами \code{title: \ldots{} | text: \ldots}; запити мають той самий пошуковий промпт. Щільний пошук із кожною збіркою йде по 4\,137 текстових документах і порівнюється з моделлю SentenceTransformers у bfloat16 на тих самих документах (\sys{dense\_eg2\_text}) з тими самими мітками. Документи, які квантований прогін приносить у першу десятку, але яких не бачив жоден суддя, вважаються нерелевантними, і їхню частку подано. Текстові колекції містять 21 рядок сторінок PDF, які повна модель кодує як зображення. GGUF натомість кодує збережений для них текст (загалом 485 символів); 84\,767 рядків мають текстовий вхід. Отже, однакові документи-кандидати не означають однакової модальності входу для кожного фрагмента, і порівняння не ізолює лише квантування.

Масштабування точного щільного пошуку виміряно поза archive-index на колекціях, вирощених із 86\,726 справжніх векторів EG2 до $10^5$, $3\cdot10^5$, $10^6$ і $2\cdot10^6$ векторів; додані вектори --- зашумлені копії справжніх (гаусів шум, $\sigma=0{,}05$, з повторним нормуванням), тож колекції синтетичні. Кожен із 200 векторів запитів UA-PAR шукається окремо оптимізованим точним вибором першої десятки в NumPy (добуток матриці й \code{argpartition}), точним індексом Faiss та індексом HNSW ($M=32$, глибина побудови 200, глибина пошуку 64 і 256). Послідовний додатковий прогін викликає справжній метод ранжування векторів archive-index на тих самих матрицях: збирання кандидатів, повне \code{argsort} і побудову 400 кортежів результатів, без перевірки версії кешу SQLite. Це відокремлює оптимізований еталон від розгорнутого алгоритму. Повнота@10 HNSW рахує збіги ідентифікаторів векторів із NumPy; однакові оцінки й різниці обчислень не зараховуються як еквівалентні ідентифікатори. Окремо перевіряються збіг точних ранжувань Faiss і NumPy та частка результатів Faiss з оцінкою не нижче десятої еталонної з допуском $10^{-6}$.

\subsection{Програмне й апаратне забезпечення}
Усі вимірювання виконано на одному ноутбуці: Intel Core i9-14900HX (32 логічні процесори), 31~ГіБ ОЗП, NVIDIA GeForce RTX 4060 Laptop GPU з 8~ГБ, Linux 7.0. Програмне забезпечення: Python 3.13.14, PyTorch 2.14.1, Transformers 5.19.0, SentenceTransformers 6.1.0, NumPy 2.5.4, SQLite 3.53.1, pymorphy3 2.0.6, ir-measures 0.4.3, llama.cpp коміт \code{8bb3cf0} і Faiss 1.15.1 в окремому середовищі для експерименту масштабування. Повний перелік збережено в \code{efficiency.json}.
